\section{What the arithmetic gap theorems give for collisions}
\label{sec:arithmetic-collisions}
The isolated-pair theorem applies directly to a certified pair of zeta
zeros. Passing from that pairwise assertion to a statistical law requires
additional information about the complementary zeros. We make the
required quantities explicit and compare them with the gap theorems.

For consecutive distinct simple ordinates $\gamma_-<\gamma_+$, put
\[
 x_\pm=2\gamma_\pm,\qquad b=\gamma_-+\gamma_+,\qquad
 \delta=2(\gamma_+-\gamma_-),\qquad
 S=\sum_{\xi\ne x_-,x_+}(b-\xi)^{-2},\quad\kappa=\delta^2S.
\]
The sum runs through the full real zero divisor of $H_0$, including
negative zeros and multiplicities. These definitions give the exact
arithmetic specialization of Theorem~\ref{thm:isolated-pair}:
for the pairs in its stated range the collision time is
\[
 a_c=\frac{(\gamma_+-\gamma_-)^2}{2}
       +\frac{(\gamma_+-\gamma_-)^4}{2}S
       +O((\gamma_+-\gamma_-)^6S^2).
\]
The relation $x=2\gamma$ accounts for both factors of two. The moving
center is $b-2a_cG'(b)/G(b)+O(\delta\kappa^{3/2})$, where
$G(z)=H_0(z)/((z-b)^2-\delta^2/4)$.

The complementary sum admits a finite-height bound. For $R\ge1$,
write
\[
 S_R=\sum_{\substack{0<|\xi-b|\le R\\\xi\ne x_-,x_+}}(b-\xi)^{-2}.
\]
The zero count in every unit interval is $O(\log(|u|+3))$, including
multiplicity. Splitting the two tails into unit intervals and comparing
with $\int_R^\infty\log(b+u+3)u^{-2}du$ gives an absolute constant
$C$ such that
\begin{equation}\label{eq:pair-exterior-tail}
 0\le S-S_R\le C\frac{\log(b+R+3)}R
 \qquad(b\ge0,\ R\ge1).
\end{equation}
For example, integration by parts bounds that integral by
$\log(b+R+3)/R+\int_R^\infty[u(b+u+3)]^{-1}du$, and the last
integral is at most $1/R$. Thus interval enclosures for finitely many
zeros, together with an explicit unit-interval zero-count constant,
produce a rigorous upper bound for $\kappa$. The gaps enter the
calculation through the denominators in $S_R$; the tail estimate does
not replace them.

The small-gap theorems recently reached half the mean spacing.
Bondarenko and Heap \cite[Theorem~1, p.~1]{bondarenko-heap-half} obtain
\[
 \liminf_{n\to\infty}
       \frac{(\gamma_{n+1}-\gamma_n)\log\gamma_n}{2\pi}\le\frac12
\]
from the zero-location theorem. Their method uses resonance for a
quotient of shifted zeta functions. Durkan and Pearce-Crump
\cite[Theorems~1--2, pp.~1--2]{durkan-pearce-crump} establish the same
lower-limit bound for every fixed primitive Dirichlet character, and
for zeta a positive proportion of gaps smaller than
$(1/2+\epsilon)$ times the mean, for every fixed $\epsilon>0$.
Their count includes repeated ordinates. Lamzouri's
\cite[Theorem~1.1, p.~2]{lamzouri-small-gaps} positive-proportion bound
at $0.7892$ is obtained without a zero-location assumption.
These statements specify the present quantitative input at a fixed
multiple of the mean spacing.

There is an elementary obstruction to deducing $\kappa\to0$ from
a bound on that multiple alone. Choose $0<\lambda<1/2$ and consider
the entire function
\[
 f_\lambda(z)=\sin(\pi z/3)\sin(\pi(z-\lambda)/3)
                         \sin(\pi(z-2\lambda)/3).
\]
Its simple real zeros have mean spacing one and repeating gap pattern
$\lambda,\lambda,3-2\lambda$. For the adjacent pair $0,\lambda$,
the third zero $2\lambda$ alone gives
\[
 \delta^2 S\ge\lambda^2(2\lambda-\lambda/2)^{-2}=4/9.
\]
The fraction of gaps of length $\lambda$ is $2/3$.
The dilation $z\mapsto z/c$ multiplies every gap by $c$ and every
term of $S$ by $c^{-2}$, so it preserves $\delta^2S$ exactly.
Thus even arbitrarily small absolute gaps and a positive fraction
of small normalized gaps can coexist with a complementary interaction
bounded below. This example disproves that proposed implication of
those two scalar gap properties; it makes no claim to reproduce the
zeta divisor.

There is also a precise limitation to using the restricted correlation
identities alone. Lagarias and Rodgers
\cite[Theorem~3.1, p.~4, with the test class in Theorem~2.4, p.~3]{lagarias-rodgers}
construct a deterministic unit-density sequence with all gaps in
$\{1/2,1,3/2,\ldots\}$ whose translation-averaged correlation
statistics agree with the sine determinant for every order and every
Schwartz test in that class. On the hyperplane
$\sum\xi_i=0$, the allowed Fourier support is
$\sum|\xi_i|<2$, exactly the restricted range used here.
Every gap in that sequence is at least $1/2$, whereas a sine process
has positive probability of a gap in $(0,\epsilon)$ for every
$\epsilon>0$. The common restricted identities therefore do not
determine these small-gap events. This counterexample concerns the
specified statistics; it does not assert that the sequence has an
Euler product or satisfies the functional equation of zeta.

The next arithmetic question is consequently specific. Can one prove
that sufficiently many adjacent \emph{distinct simple} zeta pairs have
$\delta^2\sum_{\xi\ne x_-,x_+}(b-\xi)^{-2}\to0$, with quantitative
control of nearby triples and of the remaining singleton branches?
For an arithmetic collision point-process limit, can these estimates
be combined with a marked extreme-gap limit on the shrinking scale,
including its joint law and tightness? The proved local expansion
identifies the map from such pair data to collision data, but the
fixed-scale gap results quoted above leave these estimates open.

The recent half-gap work also gives a documented example of assistance
by OpenAI's ChatGPT in neighbouring research: Bondarenko and Heap
acknowledge exploratory discussions, calculations and numerics on p.~2
of their paper. Their theorem and proof remain credited to its named
authors. The present collision estimates use their theorem only in the
comparison of the arithmetic gap input.
